% Required preamble:
% \usepackage{tikz}
% \usetikzlibrary{arrows.meta,positioning,fit,calc,backgrounds}

\begin{figure*}[t]
\centering
\resizebox{0.98\textwidth}{!}{%
\begin{tikzpicture}[
    font=\sffamily,
    >=Latex,
    node distance=7mm and 10mm,
    stage/.style={
        draw,
        rounded corners=2pt,
        line width=0.55pt,
        align=center,
        minimum width=28mm,
        minimum height=9mm,
        inner xsep=5pt,
        inner ysep=4pt,
        fill=white
    },
    module/.style={
        stage,
        minimum width=34mm,
        minimum height=11mm,
        fill=blue!4,
        draw=blue!45!black
    },
    evidence/.style={
        stage,
        minimum width=31mm,
        fill=teal!5,
        draw=teal!45!black
    },
    tool/.style={
        stage,
        minimum width=34mm,
        fill=orange!7,
        draw=orange!55!black
    },
    validator/.style={
        stage,
        minimum width=35mm,
        fill=purple!6,
        draw=purple!50!black
    },
    output/.style={
        stage,
        minimum width=31mm,
        fill=gray!8,
        draw=gray!55!black
    },
    notebook/.style={
        stage,
        minimum width=33mm,
        fill=red!5,
        draw=red!45!black
    },
    decision/.style={
        draw,
        diamond,
        aspect=1.7,
        line width=0.55pt,
        align=center,
        inner xsep=2pt,
        inner ysep=1pt,
        fill=white
    },
    flow/.style={->, line width=0.6pt, draw=black!75},
    weakflow/.style={->, line width=0.5pt, draw=black!45},
    rejectflow/.style={->, line width=0.6pt, draw=red!65!black},
    group/.style={
        draw=black!18,
        rounded corners=3pt,
        fill=black!2,
        inner sep=5pt
    },
    label/.style={font=\scriptsize\sffamily, align=center, text=black!70}
]

% Main pipeline.
\node[module] (input) {User Instruction\\+ Scenario Video};
\node[module, right=of input] (grounding) {Task-Level\\Visual Grounding};
\node[evidence, right=of grounding] (ocr) {OCR and Visual\\Evidence};
\node[module, right=of ocr] (reason) {Service-Agent\\Reasoning};
\node[decision, right=9mm of reason] (decide) {Next\\action?};

\node[tool, above right=5mm and 9mm of decide] (toolcall) {Clean JSON\\Tool Call};
\node[tool, right=of toolcall] (tools) {Official Tool\\Environment};
\node[evidence, below=7mm of tools] (toolres) {Tool Results};

\node[output, below right=5mm and 9mm of decide] (answer) {Natural-Language\\Final Answer};

\node[validator, right=16mm of decide] (valid) {Scenario-Aware\\Validator};
\node[notebook, above=7mm of valid] (notebook) {Error Notebook\\{\scriptsize retail / kitchen / restaurant / order}};
\node[decision, right=10mm of valid] (accept) {Valid?};
\node[output, right=10mm of accept] (save) {Saved Trajectory\\{\scriptsize JSON calls or final text only}};

% Internal details.
\node[evidence, below=7mm of grounding] (keyframe) {Keyframe /\\Temporal Segment};
\node[evidence, below=7mm of ocr] (bbox) {Target BBox\\and Object Region};
\node[stage, below=7mm of reason, fill=blue!2, draw=blue!25!black] (plan) {Internal Plan\\Candidate Matching\\Tool Selection};

% Group boxes.
\begin{scope}[on background layer]
    \node[group, fit=(grounding)(keyframe)(bbox), label={[label]below:Perception-side localization}] (g1) {};
    \node[group, fit=(reason)(plan), label={[label]below:Tool-grounded policy}] (g2) {};
    \node[group, fit=(valid)(notebook)(accept), label={[label]below:Experience-guided validation}] (g3) {};
\end{scope}

% Forward flow.
\draw[flow] (input) -- (grounding);
\draw[flow] (grounding) -- (ocr);
\draw[flow] (ocr) -- (reason);
\draw[flow] (reason) -- (decide);
\draw[flow] (decide) -- node[label, above left=-1pt] {tool needed} (toolcall);
\draw[flow] (toolcall) -- (tools);
\draw[flow] (tools) -- (toolres);
\draw[flow] (toolres.west) to[out=180,in=90] (reason.north);
\draw[flow] (decide) -- node[label, below left=-1pt] {task complete} (answer);

% Evidence links.
\draw[weakflow] (grounding) -- (keyframe);
\draw[weakflow] (grounding) |- (bbox);
\draw[weakflow] (keyframe) -| (ocr);
\draw[weakflow] (bbox) -- (ocr);
\draw[weakflow] (reason) -- (plan);

% Validation flow.
\draw[flow] (toolcall.east) to[out=0,in=145] (valid.north west);
\draw[flow] (answer.east) to[out=0,in=215] (valid.south west);
\draw[weakflow] (notebook) -- (valid);
\draw[flow] (valid) -- (accept);
\draw[flow] (accept) -- node[label, above] {yes} (save);
\draw[rejectflow] (accept.south) to[out=-90,in=-35] node[label, below] {high-level correction} (reason.south east);

% Small annotations.
\node[label, above=1.5mm of ocr] (ocrnote) {Noisy evidence, not final catalog facts};
\node[label, below=1.5mm of toolres] (toolnote) {Authoritative task facts};
\node[label, above=1.5mm of notebook] (nbnote) {Recurring failure patterns};

\end{tikzpicture}%
}
\caption{Overview of the proposed service-agent pipeline. The system first performs task-level visual grounding to identify the relevant temporal segment and target bounding box. OCR and visual cues are treated as noisy candidate evidence. The service agent then reasons under a strict tool-grounded policy and produces either a JSON tool call or a final natural-language response. A scenario-aware validator checks each draft using an error notebook of recurring failure patterns before the response is accepted and saved.}
\label{fig:method-pipeline}
\end{figure*}
